\appendix

\section{The Five Groups and Reference Representations}
\label{app:groups}

\begin{table}[H]
\centering
\small
\begin{tabular}{lcccl}
\toprule
$G$ & $|G|$ & $\dmin$ & solvable & realization of $\rho_G$ \\
\midrule
$S_3$ & 6   & 2 & yes & standard (dihedral) \\
$S_4$ & 24  & 3 & yes & cube rotations \\
$A_5$ & 60  & 3 & no  & icosahedral rotations \\
$S_5$ & 120 & 4 & no  & standard (zero-sum) \\
$A_6$ & 360 & 5 & no  & standard (zero-sum) \\
\bottomrule
\end{tabular}
\caption{The five task groups. $\dmin$ is the minimal faithful real
representation dimension; each reference representation is orthogonal by
construction and verified faithful (injective on $G$) at build time.
$S_4$/$A_5$ form the designed pair: matched $\dmin$, opposite
solvability.}
\label{tab:groups}
\end{table}

\section{Instrument Details and the Two Defects}
\label{app:instrument}

\textbf{Centering.} For an orthogonal target the uncentered covariance
of $Z(w)$ is isotropic and carries no subspace information; the
production lens centers it (a nontrivial irreducible representation has
zero group mean, cancelling the constant block). The step is
load-bearing: an uncentered lens scores a flawless synthetic model 0.705
where the centered lens scores 0.9996 on identical data.
% <!-- evidence: U7 -->

\textbf{Length-robustness split.} The disclosed split referenced in
Section~\ref{sec:convergence} scores only words of length $L \geq 2$
(the $L = 1$ read is attention-degenerate and was demoted with a
mechanism note in the run records): it moves per-group mean recovery
cosine by at most 0.013 (per-seed at most 0.041), with no
group-selective divergence.
% <!-- evidence: U5 -->

\textbf{The ambient-identity tax.} The first 58-cell sweep's force-rank
target was the eye-padded $\rho_G(\cdot) \oplus I_2$, whose rank is
$\dstate = \dmin + 2$ for every group.
% <!-- evidence: U6 -->
The constant identity block is
the cheapest reliable loss reduction, so a rank-capped arm buys it
first, leaving effective budget $k - 2 < \dmin$ for $\rho_G$ at every
grid point; no capped cell in that sweep ever tested the confirm
direction. The raw-artifact signature: the 39 force-rank cells' direct
cosine matched the rank-$k$ optimum $\sqrt{k/\dstate}$ with mean
absolute deviation 0.028 (maximum 0.166; the two $S_5$ below-$\dmin$
cells are the only outliers, a disclosed optimization shortfall), so the
arms had trained to their rank-constrained ceiling rather than
under-training.
% <!-- evidence: U6 -->
The sweep verdict was registered as inconclusive with the tax as
mechanism; the zero-padded causal wave of Section~\ref{sec:causal} is
the registered fix, its 30 cells configuration-verified one-by-one
against a manifest re-derived independently from the design record. In
that wave an eye-padded corroboration arm reproduces the failure on
demand at raw $k = \dmin{+}1$ (effective budget $\dmin{-}1$) while the
tax-paid point recovers.
% <!-- evidence: U3 -->

\textbf{The observational band's non-binding upper half.}
\texttt{entity\_subspace\_from\_words} returns the top-$\dmin$ left
singular vectors of the centered covariance by construction, so the
restricted state $U^{\top} Z U$ is always $\dmin \times \dmin$, and the
entropy effective rank of a $\dmin \times \dmin$ input has range
$[1, \dmin]$. Restricted effective rank therefore cannot exceed $\dmin$,
and so cannot exceed $1.3 \cdot \dmin$, for any seed of any group; only
the band's lower half, $\geq 0.7 \cdot \dmin$, is a test the data could
fail. Every group mean lands in $[0.895, 0.951] \cdot \dmin$.
% <!-- evidence: U1 -->

\textbf{The rank-constrained cosine ceiling.} The causal wave's target
is $\rho_G(\cdot) \oplus 0$, of rank exactly $\dmin$ with all $\dmin$
informative singular values equal to 1 (the reference representation is
orthogonal by construction, Appendix~\ref{app:groups}). For any matrix
$M$ of rank $\leq k \leq \dmin$, von Neumann's trace inequality gives
$\langle M, T\rangle_F \leq \sum_{i=1}^{k} \sigma_i(M)\,\sigma_i(T) =
\sum_{i=1}^{k} \sigma_i(M)$, and Cauchy--Schwarz bounds
$\sum_{i=1}^{k} \sigma_i(M) \leq \sqrt{k}\,\lVert M \rVert_F$, so
$\cos(M, T) \leq \sqrt{k/\dmin}$. At $k = \dmin{-}1$ this ceiling is
$0.707/0.816/0.816/0.866/0.894$ for $S_3/S_4/A_5/S_5/A_6$, below the
$\recninety$ threshold of 0.9 in every case: no rank-$(\dmin{-}1)$
state, however trained, can register a single word above cosine 0.9
against this target. The necessity result ($\recninety = 0.000$ at
$\dmin{-}1$, all groups, all seeds) is this bound realized, not an
independent discovery about what SGD finds; the corresponding empirical
content is that the observed cosines at $\dmin{-}1$ (0.61--0.84) sit
under their respective rank-constrained optima, evidence optimization
reaches the geometry's own ceiling. At $k \geq \dmin$ the ceiling is
1.0 (exact recovery is representable), so recovery above 0 at $\dmin$
is not guaranteed by the metric and carries the causal weight
(sufficiency).
% <!-- evidence: U3 -->

\textbf{Primary and crosscheck degauging.} The design record's default
pipeline treats scale-only degauging ($\hat{Q} = I$) as the primary
metric, with the fitted-$(c, Q)$ Procrustes score retained as a
robustness cross-check, after calibration on unconstrained checkpoints
found $\hat{Q} \approx I$ empirically. Under the force-rank grid's
zero-padded target this equivalence breaks: the informative block's
degenerate singular spectrum makes the scale-only fit's basis
arbitrary, so the pre-registration for the causal wave specifically
pins the fitted-$(c, Q)$ score, denoted $\recninety$ throughout
Section~\ref{sec:causal} and Figure~\ref{fig:razor}, as decisional; the
scale-only score informs no reported conclusion. The two diverge
sharply on this grid, which is why the distinction is disclosed here.

\section{$S_3$ Per-Seed Causal Detail}
\label{app:s3causal}

\begin{table}[H]
\centering
\small
\begin{tabular}{lcccccc}
\toprule
seed & anchor & $k{=}\dmin{-}1$ & $k{=}\dmin$ & $k{=}\dmin{+}1$ & own bar & clears \\
\midrule
0 & 0.550 & 0.000 & 0.450 & 0.550 & 0.495 & no ($-0.045$) \\
1 & 0.600 & 0.000 & 0.550 & 0.750 & 0.540 & yes ($+0.010$) \\
2 & 0.800 & 0.000 & 0.600 & 0.750 & 0.720 & no ($-0.120$) \\
3 & 0.600 & 0.000 & 0.650 & 0.600 & 0.540 & yes ($+0.110$) \\
\bottomrule
\end{tabular}
\caption{$S_3$ four-seed causal detail (crosscheck $\recninety$) behind
the seed-mean confirmation in Section~\ref{sec:causal}; the own bar is
$0.9 \times$ each seed's anchor. $k = \dmin{-}1$ reads exactly 0.000 in
all four seeds; $k = \dmin$ clears each seed's own bar in 2 of 4 (seeds
1, 3), missing in the other 2 because the anchor itself ranges
0.550--0.800 across seeds. The pre-registered criterion compares the
seed-mean $k = \dmin$ value (0.5625) against the bar fixed before the
extension ran (0.495, from seed 0's anchor), specifically to avoid
computing the bar from the same noisy seeds it is judged against;
recomputing the bar as $0.9\times$ the four-seed anchor mean (raw mean
0.6375, bar 0.574) would put the seed-mean value 0.011 below it.}
% <!-- evidence: U4 -->
\label{tab:s3seeds}
\end{table}

\section{Per-Seed Observational Rank}
\label{app:m1}

\begin{table}[H]
\centering
\small
\begin{tabular}{lccccc}
\toprule
$G$ & $\dmin$ & $n$ & restricted eff.\ rank & band & in band \\
\midrule
$S_3$ & 2 & 3 & $1.877 \pm 0.060$ & $[1.4, 2.6]$ & 3/3 \\
$S_4$ & 3 & 5 & $2.852 \pm 0.054$ & $[2.1, 3.9]$ & 5/5 \\
$A_5$ & 3 & 5 & $2.832 \pm 0.062$ & $[2.1, 3.9]$ & 5/5 \\
$S_5$ & 4 & 3 & $3.591 \pm 0.069$ & $[2.8, 5.2]$ & 3/3 \\
$A_6$ & 5 & 3 & $4.736 \pm 0.023$ & $[3.5, 6.5]$ & 3/3 \\
\bottomrule
\end{tabular}
\caption{Restricted effective rank of the unconstrained trained state
(mean $\pm$ sd over seeds) against each group's $\dmin$, with the
pre-registered $[0.7, 1.3] \cdot \dmin$ band. Every seed of every group
is in band; Spearman $\rho = 0.9747$ is the tie-capped maximum; the
dimension-matched pair $S_4$/$A_5$ differs by 0.019 rank-units.}
% <!-- evidence: U1 -->
\label{tab:m1}
\end{table}
